% Specification review / feedback for SQIsign authors
% Derived from implementing sqisign-selkie (Rust, NIST-I) against v3.0
\documentclass[a4paper,11pt]{article}

\usepackage{amsfonts,amsmath,amssymb}
\usepackage[T1]{fontenc}
\usepackage{libertinus}
\let\arrowvert\undefined
\let\Arrowvert\undefined
\usepackage[libertine]{newtxmath}
\usepackage{microtype}
\usepackage[dvipsnames]{xcolor}
\usepackage{enumitem}
\usepackage{booktabs}
\usepackage{longtable}
\usepackage{array,fancyvrb,needspace}
\usepackage[colorlinks=true, urlcolor=BrickRed, linkcolor=BrickRed,
            citecolor=BrickRed]{hyperref}
\usepackage{xurl}
\usepackage{geometry}
\geometry{margin=28mm, marginparwidth=0pt}
\setlength{\emergencystretch}{3em}

\newcommand{\algo}[1]{Algorithm~#1}
\newcommand{\Fp}{\mathbb{F}_p}
\newcommand{\Fpii}{\mathbb{F}_{p^2}}
\newcommand{\Oo}{\mathcal{O}_0}
\newcommand{\code}[1]{\texttt{#1}}
\newcommand{\file}[1]{\texttt{#1}}
\newcommand{\nrd}{\operatorname{nrd}}

% Security severity and implementation-priority badges.
\newcommand{\blocking}{\textsc{\color{BrickRed}Blocking}}
\newcommand{\high}{\textsc{\color{BurntOrange}High}}
\newcommand{\medium}{\textsc{\color{RoyalBlue}Medium}}
\newcommand{\low}{\textsc{\color{OliveGreen}Low}}

% Finding header: affected algorithms + severity on one line
% Badge flush right on the same line when it fits, else flush right on
% the next line.
\newcommand{\finding}[2]{\noindent\textit{#1}%
  \unskip\nobreak\hfil\penalty50\hskip1em\hbox{}\nobreak\hfill#2%
  {\parfillskip=0pt\finalhyphendemerits=0\par}\smallskip}


\newcommand{\interopfinding}[2]{\finding{#1}{\textit{Conformance:} #2}}
\newcommand{\Hsim}{H_{\mathrm{sim}}}
\newcommand{\Hunif}{H_{\mathrm{unif}}}
\newcommand{\Drsp}{D_{\mathrm{rsp}}}
\newcommand{\specurl}{https://sqisign.org/spec/sqisign-20260901.pdf}
\newcommand{\specat}[2]{\href{\specurl\#nameddest=#1}{#2}}
\renewcommand{\arraystretch}{1.12}
\hypersetup{pdftitle={SQIsign v3: Specification Review},
            pdfauthor={Deirdre Connolly, Selkie Cryptography},
            pdfsubject={Findings and suggested changes for the SQIsign team},
            linktoc=all,bookmarksopen=true,bookmarksnumbered=true}

\title{SQIsign v3: Specification Review}
\author{Deirdre Connolly\\[0.2em]
  \small Selkie Cryptography\\
  \small\texttt{durumcrustulum@gmail.com}}
\date{7 September 2026\enspace{\small(draft for discussion)}}

\AddToHook{cmd/section/before}{\Needspace{8\baselineskip}}
\AddToHook{cmd/part/before}{\Needspace{18\baselineskip}}

\begin{document}
\maketitle

\noindent Notes on the 1 September 2026 v3 specification~\cite{sqisign-spec}
and the round-3 reference implementation~\cite{sqisign-cref}, with
suggested changes for the SQIsign team.

The interval-sampling defect is the main confirmed security finding.
The hint-distribution concern remains provisional. Neither establishes
a signature forgery or key-recovery attack. Part~\ref{part:security}
sets out those findings and the related deployment considerations.
Parts~\ref{part:conformance} and~\ref{part:clarifications} collect the
implementation corrections, ordered by conformance priority:
\emph{blocking} for fundamental algorithm or random-stream mismatches,
then high, medium, and low. These priorities describe implementation
impact, separately from the security severity in Part~\ref{part:security}.
Part~\ref{sec:v2} records the 42 v2 findings.

\begingroup
\small
\tableofcontents
\endgroup
\clearpage

\part{Security findings and deployment}
\label{part:security}

\section{Biased interval sampling}
\label{sec:sampler}
\hypertarget{SQI3-SPEC-01}{}
\finding{\algo{4.24} and Section 4.1.4}{\medium}

\specat{func.SampleFromInterval}{\textsc{SampleFromInterval}}
sets $m=b-a+1$ and $n=\code{ibz\_bitsize}(m)$, draws an $n$-bit integer $x$,
and accepts it when
\[
 x<m\left\lceil\frac{2^n}{m}\right\rceil.
\]
The threshold is at least $2^n$, so every draw is accepted. Reducing $x$
modulo $m$ therefore biases the output whenever $m$ is not a power of two.
For $[a,b]=[0,2]$, the equiprobable draws $0,1,2,3$ produce $0,1,2,0$:
the output probabilities are $1/2,1/4,1/4$, rather than $1/3$ each.

More generally, write $R=2^n=qm+r$, with $0\leq r<m$. The first $r$
residues have $q+1$ preimages and the others have $q$. Their total variation
distance from uniform is
\[
 \Delta=\frac{r(m-r)}{mR}.
\]
The prose describes 64 extra random bits, which line~2 omits. Extra bits
would reduce this bias, but would not fix the acceptance test.

For the \code{EIBox} values in Table~9, the coefficient draws used by
\algo{4.78} have the following distributions before later rejections:

\begin{center}
\begin{tabular}{@{}llll@{}}
\toprule
Level & Interval & Distance from uniform & Mean (uniform: zero)\\
\midrule
I & $[-15,15]$ & $15/496\approx 3.02\%$ & $-15/32$\\
III & $[-19,19]$ & $175/1248\approx 14.02\%$ & $-175/64$\\
V & $[-22,22]$ & $247/1440\approx 17.15\%$ & $-247/64$\\
\bottomrule
\end{tabular}
\end{center}

These are biases in internal draws, not measured biases in public keys or
signatures. They do not describe the C reference's wider sampler either.

\Needspace{8\baselineskip}
\paragraph{Why this matters to the analysis.}
The response path is \textsc{Sign} $\to$ \textsc{SampleQuaternionResponse}
$\to$ \textsc{RandomBoundedElement} $\to$ \textsc{SampleFromInterval}.
The sampler also feeds prime-norm ideal sampling and auxiliary ideal
sampling through Algorithms~4.71 and~4.82.

Uniform coordinate sampling in a box enclosing a lattice ball, followed by
rejection to the ball and to odd norm, gives equal weight to the accepted
lattice points. Biased coordinate sampling loses that argument. The weights
can depend on the chosen basis, whose representation can depend on secret
ideal data. I have not established whether that dependence is visible in
complete signatures; correcting the sampler restores the intended starting
point for the distribution argument.

The accompanying script also gives a small $\mathbb{Z}^4$ example where
odd-norm rejection and identifying $v$ with $-v$ leave unequal weights.
Those operations are therefore not a general cure for the bias, although
this example is not a SQIsign response lattice.

\Needspace{12\baselineskip}
\paragraph{Suggested change.}
Use an acceptance range whose size is divisible by $m$: replace the
ceiling with a floor. An exact sampler is:

\begin{Verbatim}[fontsize=\small,frame=single,framesep=3mm]
m = b - a + 1
if m == 1: return a
n = bit_length(m) + kappa
R = 2^n
limit = R - (R mod m)
repeat:
    x = uniform_n_bit_integer(prng)
until x < limit
return a + (x mod m)
\end{Verbatim}

Here $\kappa\geq0$ changes the rejection rate, not the uniformity of accepted
outputs. If approximately uniform fixed-draw sampling is preferred, I
suggest specifying its statistical-distance budget and accounting for all
calls and subsequent conditioning in the analysis. The prose, pseudocode,
and reference should agree on which contract is intended.

This pseudocode specifies the distribution only. Public working-width
bounds, observable retries, and failure handling need a separate timing
argument. A fixed candidate batch with masked selection is one possible
approach, provided its all-rejected probability and failure behavior are
accounted for.

\paragraph{Reference implementation and conformance.}
The C reference uses a tracked bit bound plus 64 bits~\cite{cref-sampler}.
It adds $m$ to
the draw, checks that the sum still fits, and then reduces modulo $m$.
In a release build the accepted sample space has size $R-m$; with assertions
enabled the exceptional branch asserts. Since $R-m$ is not generally a
multiple of $m$, this is not exact uniform sampling either, though the bias
is much smaller than in the printed algorithm.

Our current Rust interval sampler already uses exact rejection over
$[0,b-a]$. I suggest preserving that distribution during the port, while
addressing its variable-time behavior separately. A changed random-byte
schedule may require new KATs without changing the wire format or
verification relation. Useful conformance cases include singleton and
signed intervals, powers of two and adjacent widths, rejection boundaries,
and the three \code{EIBox} intervals.

\Needspace{8\baselineskip}
\section{Hint degree weighting}
\label{sec:proof}
\hypertarget{SQI3-PROOF-02}{}
\finding{Figure 4; Problem 8.1.5 and Section 8.2.2.2}{\low\ (provisional)}

In \specat{fig:sqihintdistrs}{Figure~4}, $\Hunif$ samples an odd degree
$d\leq\Drsp$ with weight
\[
 \psi(d)=\prod_{\ell^e\parallel d}(\ell+1)\ell^{e-1},
\]
then samples a degree-$d$ isogeny with non-cyclic kernels allowed. For $d$
coprime to $p$, $\psi(d)$ counts cyclic kernels. Counting all degree-$d$
kernels instead gives
\[
 \sigma_1(d)=\sum_{a\mid d}a
 =\prod_{\ell^e\parallel d}(1+\ell+\cdots+\ell^e).
\]
Here isogenies are counted up to isomorphism of the target, without
quotienting by automorphisms of the source. Section~5.6 of the earlier
security proof also identifies $\psi$ as the cyclic count~\cite{proof}.

For example, degree~9 has twelve cyclic kernels and one additional kernel,
$E[3]$, for a total of thirteen. Degree~25 similarly gives 30 versus 31.
The printed procedure assigns an individual degree-$d$ isogeny probability
proportional to $\psi(d)/\sigma_1(d)$, which varies with $d$. Thus it is
not uniform over all outgoing bounded odd-degree isogenies. This is a
well-defined distribution, but its degree dependence needs to be reconciled
with the intended comparison to $\Hsim$.

The latter first chooses a target from the stationary distribution, then
samples uniformly among bounded odd-degree isogenies to that target.
\textsc{RandomBoundedElement} does not impose primitivity or remove odd
scalar factors, so we cannot simply assume that every response is cyclic.

\paragraph{A statistic worth checking.}
Consider $3\mid\deg(\varphi_2)$. If the response homomorphisms are uniform
modulo~3 in $M_2(\mathbb{F}_3)$, this event has probability
\[
 \frac{\#\{M:\det M=0\}}{3^4}=\frac{33}{81}=\frac{11}{27}.
\]
The $\psi$-weighted degree distribution, over increasing odd-degree bounds,
instead gives limiting probability $2/5$. The gap is $1/135$. The event
can be checked from the degree, or from the rank of the action on 3-torsion
when the hint representation makes that action available.

The missing step is a quantitative residue-distribution argument for
$\Hsim$ at the actual joint growth of $p$ and $\Drsp$, including target
conditioning and exceptional curves. The script checks the kernel counts,
the 81 small-field matrices, and finite degree-weight marginals. It does
not simulate $\Hsim$ or prove this distinguisher.

\paragraph{Suggested follow-up.}
If the intended measure is uniform over all outgoing bounded odd-degree
isogenies, I suggest investigating $\sigma_1(d)$ as the degree weight and
rechecking the resulting proof distribution and its transport properties.
If the $\psi$ weights are intentional, an explanation of the degree
dependence and the statistic above would help settle the question.

Theorem~8.1.7 already includes a hint-distinguishing advantage term. This
observation therefore questions whether the justification makes that term
small; it does not contradict the reduction as written. Nor does it refute
the separate argument using the hardness of the $\Hsim$-assisted relation
directly. I would treat this initially as a change or clarification to the
analysis. Forcing primitive or cyclic responses in the signer would be a
different protocol change requiring its own analysis.

\Needspace{8\baselineskip}
\section{Deployment properties to specify}
\label{sec:related}

The documented timing and binding limitations below affect deployment
choices. The verifier and work-bound requests connect to the detailed
implementation findings later in this review.

\paragraph{Timing resistance.}
Section~8.5 already explains that signing and key generation are not fully
constant-time. For a hardened version, it would help to state which widths,
iteration bounds, and failure behavior are public, and which operations
remain outside the timing claim. In particular, correcting interval
sampling does not make \textsc{IdealToIsogeny} or the whole signer
constant-time. This needs resolution before deployment wherever the
adversary can observe the relevant secret-dependent behavior.

\paragraph{Stronger message binding.}
The optional construction in Section~8.4.1 would benefit from explicit
signing, verification, and encoding rules: carry the full $2\lambda$-bit
digest, use the designated $\lambda$ bits for the challenge walk, and
\emph{compare all $2\lambda$ bits} during verification. Checking only the
walk bits would discard the extra binding. This adds 16, 24, or 32 bytes
at levels I, III, or V. Any incompatible profile needs a distinct format
identifier and test vectors, including rejection when only the extra
digest bits are changed.

As Section~8.4.2 explains, the auxiliary-isogeny representation remains
malleable. The extended digest does not provide strong unforgeability,
and matrix normalization does not make complete signatures unique.

\paragraph{Verifier preconditions and failure behavior.}
Section~3.1.2 already requires $A^2\ne4$ for \code{ECCurve}, and the C
verifier checks the public and auxiliary coefficients before reconstructing
bases~\cite{cref-verify}. I suggest making that checked decoding boundary
explicit in \algo{3.24}, along with propagation of supporting-routine
failures to rejection. Similarly, map each deferred torsion, independence,
kernel, and splitting check to the operation that enforces it. The inspected
gluing code does check torsion and independence~\cite{cref-gluing}; I have
not established the suspected missing odd-response check as a bypass.

Sections~\ref{sec:verify} and~\ref{sec:chainchecks} detail the verifier
clarifications; \S\ref{sec:fromhint} covers a justified work bound and
failure policy for zero-hint basis reconstruction. Negative vectors for
malformed encodings, singular coefficients, invalid kernel data, and
failure propagation would make these requirements concrete.

\part{Implementation conformance}
\label{part:conformance}

\section{Incorrect norm equation in LoopTwo}
\label{sec:looptwo}

\interopfinding{Affects \algo{4.148} (\textsc{LoopTwo}) and \algo{4.149}
  (\textsc{IdealToIsogenyNormEquation}), hence every
  \textsc{IdealToIsogeny} call.}{\blocking}

Line~10 sets $d \leftarrow 2^e - 2\mu^2 r_0 - \mu N$.  The comment on
line~11 says the target is $2^e - 2\mu^2 r_0 - \delta N$, and $\delta$
(lines~2 to~4) is otherwise unused.  The reference
(\code{quat\_\allowbreak{}qlapoty\_\allowbreak{}loop\_\allowbreak{}two}, \code{quaternion/\allowbreak{}ref/\allowbreak{}lvlx/\allowbreak{}qlapoty.c})
uses $\delta N$ with $\delta = 3$ when $r_0$ is even and $1$
otherwise.  With $\mu N$ the value $z$ of line~16 is not the one the
Cornacchia step needs, and the returned $\beta_1, \beta_2$ do not
satisfy $\nrd(\beta_1) + \nrd(\beta_2) = N 2^e$.

Two further differences on the failure path:
\begin{itemize}[nosep]
  \item The reference iterates while
    $\mu \le \lfloor \sqrt{2^e / r_0} \rfloor$. Before rounding and
    endpoint conventions, this upper bound is twice that of line~6
    ($\mu < \sqrt{2^{e-2}/r_0}$).
  \item On \textsc{LoopTwo} failure \algo{4.149} line~13 reruns
    \textsc{LoopOne}.  The reference (\code{quat\_\allowbreak{}qlapoty\_\allowbreak{}normeq})
    reports failure and its caller
    (\code{dim2id2iso\_\allowbreak{}ideal\_\allowbreak{}to\_\allowbreak{}isogeny\_\allowbreak{}qlapoty},
    \code{id2iso/\allowbreak{}ref/\allowbreak{}lvlx/\allowbreak{}dim2id2iso.c}) aborts.
\end{itemize}

\paragraph{Recommendation.}
Write $\delta N$ on line~10.  State the loop bound the reference uses
and whether a rerun of \textsc{LoopOne} is permitted.

\section{Sampling gates and ideal canonicalization}
\label{sec:smallest}

\interopfinding{Affects \algo{4.80} (\textsc{SmallestEquivalentIdeal}) via
  \algo{4.149} line~1, i.e.\ every \textsc{IdealToIsogeny}
  call.}{\blocking}

The reference (\code{quat\_\allowbreak{}ideal\_\allowbreak{}shortest\_\allowbreak{}equivalent},
\code{quaternion/\allowbreak{}ref/\allowbreak{}lvlx/\allowbreak{}ideal.c}) computes $\alpha_1$ as the first
vector of the reduced basis and $N' = \nrd(\alpha_1)/N$, then:
\begin{enumerate}[nosep]
  \item It draws $\beta$ only if $\gcd(N', N) \ne 1$.  \algo{4.80}
    always enters the sampling loop, and its loop condition is
    $\gcd(N', \nrd(\beta)/N) = 1$, a different test.
  \item The generator is $\overline{\alpha_1}$, or
    $\beta \cdot \overline{\alpha_1} / N$ after a draw.  Line~9 uses
    $\alpha_1$ without conjugation.
  \item After building $(N', (x, y))$ it canonicalizes: if $2x > N'$
    then $(x, y) \leftarrow (N' - x,\, N' - y - 1)$ and the returned
    generator is multiplied by $\mathbf{i}$ on the right.  \algo{4.80}
    has no such step.
  \item Line~1 calls \textsc{InertIdealReduction} without the bound
    $\beta_\lambda$ that \algo{4.72} requires.
\end{enumerate}
Item~1 means a literal \algo{4.80} consumes default-domain bytes on
every call, shifting every later draw from that stream (\algo{4.23}
witnesses, \algo{4.146}).  Item~3 replaces $J$ by $J\mathbf{i}$, which
has the same codomain but a different image basis, so
$M_{\mathrm{sk}}$ and $M_{\mathrm{chl}}$ can differ when that
normalization changes the representative.

\paragraph{Recommendation.}
Specify the gated draw, the conjugate, and the $2x \le N'$ canonical
form. Explain which output representatives this normalization fixes,
including the pushed basis; reproducing those choices is necessary
for byte-for-byte conformance.

\section{Random-byte consumption is underspecified}
\label{sec:sample}
\interopfinding{\algo{4.24} and callers in ideal and response sampling}{\blocking}

Section~4.1.4 describes a draw of $\code{ibz\_bitsize}(b-a+1)+64$ bits,
while \algo{4.24} draws only $\code{ibz\_bitsize}(b-a+1)$ bits. The
distributional defect and the exact rejection correction are discussed in
\S\ref{sec:sampler}. A separate conformance issue is the choice of working
width and the resulting random-byte schedule.

The reference's \code{ibz\_rand\_interval\_with\_domain} draws
$\lceil(\code{b->bitlen}+64)/8\rceil$ bytes and masks to
$\code{b->bitlen}+64$ bits~\cite{cref-sampler}. Here \code{bitlen} is
the tracked bound of Definition~4.1.1, rather than the exact bit size of
the endpoint or interval width. The reference then adds the interval
width, checks that the sum fits, reduces, and shifts by $a$; its
overflow check is not the exact rejection rule recommended above.
The spec does not give the tracked bound at each call site.

The earlier implementation comparison identified these byte schedules:
\begin{itemize}[nosep]
  \item The sign draw in \algo{4.71} uses $[0,1]$ with the bound of the
    constant one, consuming nine bytes per candidate.
  \item \algo{4.83} samples $[0,2\beta_i]$ and shifts by $-\beta_i$;
    its working width comes from the bound on $2\beta_i$
    (\code{quat\_lattice\_sample\_from\_ball},
    \path{quaternion/ref/lvlx/lattice.c}).
  \item \algo{4.78} uses
    \code{ibz\_rand\_interval\_minm\_m}, which samples $[0,2m]$ with a
    33-bit bound and subtracts $m$. For $m=15$, each coefficient draw
    consumes thirteen bytes.
  \item \algo{4.46} assigns $\lceil\log_2 p\rceil$ as the endpoint bound.
\end{itemize}
Even a one-byte sign draw instead of nine bytes changes subsequent draws
from that stream. A mathematically equivalent sampling interface therefore
need not reproduce the reference KATs.

\paragraph{Suggested change.}
Specify the intended distribution first, then the public working-width
bounds, PRNG domain, byte order, masking, singleton handling, and retry
behavior needed for reproducibility. Tabulate the bounds at each call
site, including the signed-interval variant. Do not substitute the exact
size of a secret-dependent endpoint for its public bound without a timing
argument. Updating the sampler and KATs can preserve the existing wire
format and verification relation.

\section{Incorrect product-ideal basis}
\label{sec:mulideal}

\interopfinding{Affects \algo{4.73} (\textsc{MulIdealBasis}) and \algo{4.75}
  (\textsc{MulIdealDualReduction}) line~2, hence \algo{4.83} and every
  response.}{\high}

\algo{4.74} builds $Y_1 = N_2(-u_2(2y_1+1) \bmod 2N_1)$ and
$Y_2 = N_1(u_1(2y_2+1) \bmod 2N_2)$.  Reducing modulo $2N_i$ preserves
parity, so $Y_1 \equiv N_2 u_2$ and $Y_2 \equiv N_1 u_1 \pmod 2$.  With
$u_1 N_1 + u_2 N_2 = 1$ and $N_1, N_2$ odd this gives
$Y_1 - Y_2 \equiv Y_1 + Y_2 \equiv 1 \pmod 2$.

\algo{4.73} prints the third column as
$X_1 + X_2 + \tfrac{Y_1 - Y_2}{2}\mathbf{i} + \tfrac{\mathbf{i}+\mathbf{j}}{2}$.
Its $\mathbf{j}$-coefficient is $\tfrac12$, so membership in $\Oo$
requires a half-integer $\mathbf{i}$-coefficient.  The printed
coefficient $\tfrac{Y_1 - Y_2 + 1}{2}$ is an integer.  The fourth
column fails the same way: its $\mathbf{k}$-coefficient is $\tfrac12$
but its constant term $\tfrac{1 - Y_1 - Y_2}{2}$ is an integer.
Neither column lies in $\Oo$, so the printed lattice is not the
product ideal.

The reference (\code{quat\_\allowbreak{}ideal\_\allowbreak{}mul\_\allowbreak{}O0}) uses
\[
  X_1 + X_2 + \tfrac{Y_2 - Y_1}{2}\mathbf{i} + \tfrac{\mathbf{j}}{2},
  \qquad
  -\tfrac{Y_1 + Y_2}{2} + (X_2 - X_1)\mathbf{i} + \tfrac{\mathbf{k}}{2},
\]
whose coefficients have the right parities.  Equivalently, \algo{4.75}
line~2 should read
$A, B, C, E \leftarrow 2(X_1 + X_2),\; Y_2 - Y_1,\; -(Y_1 + Y_2),\; 2(X_2 - X_1)$.
We confirmed on small primes that the reference's lattice equals
$\overline{I_1} \cdot I_2$, as \algo{4.74} promises, and that the
printed one does not.

\paragraph{Recommendation.}
Correct the two columns in \algo{4.73} and line~2 of \algo{4.75}.

\section{Missing response basis point and reversed outputs}
\label{sec:hdresponse}

\interopfinding{Affects \algo{3.21} (\textsc{HDResponse}) and \algo{3.23}
  lines~17 and~18.}{\high}

\begin{enumerate}[nosep]
  \item $e$ is used on lines~9 to~12 and never bound.  The output
    types imply $e = e_{\mathrm{rsp}}$.
  \item Line~12 pushes $(P_{\mathrm{com}}, 0)$ and
    $(Q_{\mathrm{com}}, 0)$ only, but the outputs
    $\mathcal{B}_{\mathrm{aux}}, \mathcal{B}_{\mathrm{chl}}$ are typed
    as bases and \algo{3.23} line~18 feeds them to
    \textsc{SetChangeOfBasisMatrix}, whose \textsc{ChangeOfBasis}
    calls need $R - S$.  The reference
    (\code{compute\_\allowbreak{}dim2\_\allowbreak{}isogeny\_\allowbreak{}challenge},
    \code{signature/\allowbreak{}ref/\allowbreak{}lvlx/\allowbreak{}sign.c}) pushes $P$, $Q$, and $P - Q$.
    Recomputing the difference on the codomain is not an option, see
    Part~\ref{sec:v2}, item~17.
  \item By the convention of \algo{3.20} ($\deg \varphi_1 = q$ from
    $E_1$), the first codomain factor for $E_1 = E_{\mathrm{com}}$ is
    $E_{\mathrm{chl}}$ and the second is $E_{\mathrm{aux}}$.  Line~12
    destructures them as $E_{\mathrm{aux}}, E_{\mathrm{chl}}$.  The
    reference reads $E_{\mathrm{chl}}$ from the first factor and
    $E_{\mathrm{aux}}$ from the second.  A signer following line~12
    encodes the challenge curve as $E_{\mathrm{aux}}$ and fails its
    own verifier.
  \item Line~7 inverts $\mathrm{IdealNorm}(I_{\mathrm{sk,chl}})$
    modulo $2^f$; the reference inverts modulo $2^{e_{\mathrm{rsp}}+2}$.
    The results agree after the reduction on lines~10 and~11.
  \item Line~12 writes $0_{E_{\mathrm{com}}}$ for a point of
    $E'_{\mathrm{aux}}$.
\end{enumerate}

\paragraph{Recommendation.}
Bind $e$, push the third point, and swap the two curves on line~12.

\section{Incorrect power-of-two ideal formulas}
\label{sec:powtwo}

\interopfinding{Affects \algo{4.68} (\textsc{InertGenToIdealPowTwo}) via
  \algo{4.151} line~14, i.e.\ the challenge ideal of every
  signature.}{\high}

Line~10 reads $X \leftarrow t(ys + zy + yz) \bmod 2^e$.  With $s$ the
inverse computed on line~9 this is not a formula for $X$.  The
reference (\code{quat\_\allowbreak{}ideal\_\allowbreak{}create\_\allowbreak{}O0\_\allowbreak{}pow\_\allowbreak{}two},
\code{quaternion/\allowbreak{}ref/\allowbreak{}lvlx/\allowbreak{}ideal.c}) computes, with
$(t, x, y, z) = \textsc{CoordinatesO}_0(\gamma \bmod 2^e)$ and
$s = (y^2 + z^2)^{-1} \bmod 2^e$,
\[
  X = s\,(ty + xz + yz) \bmod 2^e, \qquad
  Y = s\,(xy - tz - z^2) \bmod 2^e .
\]
The reference also normalizes parities before these formulas: after
the $(\mathbf{i}+\mathbf{j})/2$ step of lines~3 to~5 it multiplies
$\gamma$ by $\mathbf{i}$ on the left whenever $y$ is odd, so that the
formulas see $y$ even and $z$ odd.  \algo{4.68} accepts
$(y, z) \bmod 2 \in \{(0,1), (1,0)\}$ and applies one formula to both.

\paragraph{Recommendation.}
Replace lines~10 and~11 with the formulas above and add the
$\mathbf{i}$ normalization, or give formulas valid for both parity
cases.

\section{Inconsistent challenge-kernel coordinates}
\label{sec:kerdecomp}

\interopfinding{Affects \algo{3.16}, \algo{4.151}
  (\textsc{KernelDecomposedToIdeal}), and \algo{3.23} lines~12
  and~13.}{\high}

Three statements about the input $(c_1, c_2)$ disagree:
\begin{itemize}[nosep]
  \item \algo{3.16}: both divisible by $2^{f-e}$, not both by
    $2^{f-e+1}$, and $E_0[I] = \langle [c_1]P_0 + [c_2]Q_0 \rangle$ on
    the $2^f$-torsion basis $\mathcal{B}_0$.
  \item \algo{4.151}: ``a point $[c_1]P_0 + [c_2]Q_0$ of order $2^e$''.
  \item \algo{3.23} line~12: $(c_1, c_2) = M_{\mathrm{sk}} \cdot
    [1, a_{\mathrm{chl}}]^T$ with $M_{\mathrm{sk}}$ reduced modulo
    $2^\lambda$, which are coordinates on the $2^\lambda$-torsion
    basis $[2^{f-\lambda}]\mathcal{B}_0$ and are not divisible by
    $2^{f-\lambda}$.
\end{itemize}
The reference (\code{id2iso\_\allowbreak{}kernel\_\allowbreak{}dlogs\_\allowbreak{}to\_\allowbreak{}ideal\_\allowbreak{}even},
\code{id2iso/\allowbreak{}ref/\allowbreak{}lvlx/\allowbreak{}id2iso.c}) takes the third reading: it evaluates
$M_{\mathrm{sk}}$ on $(1, a_{\mathrm{chl}})$ and reduces every
intermediate modulo $2^{\lambda}$.  An implementor who follows
\algo{3.16} multiplies by $2^{f-\lambda}$ first and computes the
ideal of a different point.  \algo{4.151} line~1 also uses
$M_{\mathbf{j}}$, which is not among the three precomputed matrices of
\S4.6; it equals $2M_{(\mathbf{i}+\mathbf{j})/2} - M_{\mathbf{i}}$.

\paragraph{Recommendation.}
State once that $(c_1, c_2)$ are coordinates modulo $2^e$ on
$[2^{f-e}]\mathcal{B}_0$, and fix \algo{3.16} to match.

\section{Hint exponents and the base-curve basis}
\label{sec:hints}

\interopfinding{Affects \algo{3.11} (\textsc{TorsionBasisToHint}),
  \algo{3.9} lines~1 and~2, \algo{4.94}, \algo{4.95}.}{\high}

\algo{3.11} says the hint decodes to the basis
$([2^{e_h}]P, [2^{e_h}]Q)$ of $E[2^{e - e_h}]$, and \algo{4.94} types
the output as $\mathrm{Hint}_{e - e_h}$.  \algo{3.9} calls
$\textsc{TorsionBasisToHint}(E_1, f, e)$ with $e = e_{\mathrm{rsp}}+2$,
which under \algo{3.11} would produce a hint for $E[2^{f - e}] =
E[2^{126}]$.  \algo{3.24} line~9 decodes the same hint at exponent
$e_{\mathrm{rsp}} + 2$.

\algo{4.94} lines~21 to~26 resolve it: $e_h$ is the exponent at which
\textsc{TorsionBasisFromHint} will be called, and the adjustment
selects between $[2^{e-e_h}](R \pm S)$ so that the third coordinate
matches the one \algo{4.95} recomputes.  The reference
(\code{ec\_\allowbreak{}curve\_\allowbreak{}to\_\allowbreak{}basis\_\allowbreak{}2f\_\allowbreak{}to\_\allowbreak{}hint}, \code{ec/\allowbreak{}ref/\allowbreak{}lvlx/\allowbreak{}basis.c})
agrees.  Only the prose of \algo{3.11} and the type in \algo{4.94} are
wrong.

A second gap: when $A=0$ the reference returns the precomputed basis
of $E_0$ with hint~0, and \textsc{FromHint} likewise handles this curve
separately. Algorithm~4.94 has no such branch. Its candidate has
$x_R=0$ for every $h$; $(0,0)$ is a point of order two on $E_0$, and
the construction gives $S=R$ before cofactor multiplication. It cannot
produce the required basis. Hint~0 therefore covers both the base
curve and a fallback search; the reference distinguishes them by
testing $A$.

\paragraph{Recommendation.}
Rewrite \algo{3.11} as ``return a basis of $E[2^e]$ and a hint that
\textsc{TorsionBasisFromHint}$(E, e_h, \cdot)$ decodes to a basis of
$E[2^{e_h}]$ with the same third coordinate''.  Add the $A = 0$
branch to \algo{4.94} and \algo{4.95}.

\section{Wrong symmetric element in coordinate conversion}
\label{sec:thetachange}

\interopfinding{Affects \algo{4.119} (\textsc{ThetaChangeOfCoordinates})
  lines~14 to~21, hence every gluing.}{\high}

Line~3 sets $(\mathbf{G}^1, \mathbf{G}^2, \mathbf{G}^3)$ to the
symmetric elements of $K'_{1,1}$, $K'_{2,1}$, and their sum on $E_1$.
Lines~14 to~21 build rows~2 and~3 of $\mathbf{M}$ from
$\mathbf{G}^2$, the element of $K'_{2,1}$.  The reference
(\code{gluing\_\allowbreak{}change\_\allowbreak{}of\_\allowbreak{}coordinates}, \code{hd/\allowbreak{}ref/\allowbreak{}lvlx/\allowbreak{}gluing.c})
builds those rows from \code{Basis\_\allowbreak{}1.g1}, the element of
$K'_{1,1}$, that is $\mathbf{G}^1$.  Row~1 uses $\mathbf{H}^2$ in both.
The index pattern of the products is otherwise identical.

\paragraph{Recommendation.}
Replace $\mathbf{G}^2$ by $\mathbf{G}^1$ in lines~14 to~21, or explain
why the two agree.

\section{Conjugation and splitting in coprime ideal sampling}
\label{sec:coprime}

\interopfinding{Affects \algo{4.79} (\textsc{EquivalentCoprimeIdeal}) via
  \algo{4.81} (both keys), \algo{3.18} line~1 (every
  signature).}{\high}

The reference (\code{quat\_\allowbreak{}ideal\_\allowbreak{}small\_\allowbreak{}equivalent\_\allowbreak{}coprime\_\allowbreak{}enumeration},
\code{quaternion/\allowbreak{}ref/\allowbreak{}lvlx/\allowbreak{}ideal.c}) differs from \algo{4.79}:
\begin{enumerate}[nosep]
  \item Coefficients come from the $[-m, m]$ sampler variant
    (\S\ref{sec:sample}), 13 bytes each, and an all-zero vector is
    skipped before the content is divided out.
  \item $N'$ must be odd before the primality or coprimality test.
    For $M = 2$ oddness is the whole test.
  \item $\beta$ is conjugated before the ideal is built.
  \item A non-inert $\chi_I(\beta)$ is not rejected.  The reference
    calls \code{quat\_\allowbreak{}ideal\_\allowbreak{}create\_\allowbreak{}O0\_\allowbreak{}odd}, which splits off the
    $\mathbb{Z}[\mathbf{i}]$ factor as in \algo{4.76} and returns the
    inert part.  \algo{4.79} line~10 resamples instead.
  \item The loop has no $\mathrm{EIBoundIteration}$ cap.
\end{enumerate}
Item~3 changes which ideal is produced from a given $\beta$ and
item~4 changes how many draws are made before returning.  Both alter
$I_{\mathrm{sk}}$ and $I_{\mathrm{com}}$ on inputs that occur in the
published KATs.  Line~9 also passes the arguments of
\textsc{InertGenToIdeal} in the opposite order from \algo{4.70}.

\paragraph{Recommendation.}
Reconcile the conjugation, oddness test, and non-inert splitting with
the intended ideal distribution, then state one algorithm for
implementations to follow. Specify whether the iteration bound is
normative and how its failure policy relates to the analysis.

\section{Response draw skipping and iteration policy}
\label{sec:rbe}

\interopfinding{Affects \algo{4.83} (\textsc{RandomBoundedElement}) via
  \algo{3.18} line~5.}{\high}

The reference (\code{quat\_\allowbreak{}lattice\_\allowbreak{}sample\_\allowbreak{}from\_\allowbreak{}ball}) differs from
\algo{4.83} in its random-byte schedule and failure behavior:
\begin{enumerate}[nosep]
  \item When $\beta_i = 0$ it sets $x_i = 0$ without consuming
    randomness.  \algo{4.83} line~7 always draws.
  \item Each coordinate is drawn from $[0, 2\beta_i]$ and shifted, see
    \S\ref{sec:sample}.
  \item The loop is unbounded; there is no
    $\mathrm{BoundIterationRBE}$ check.
\end{enumerate}
The acceptance test is equivalent: the reference accepts when
$x^t G x \not\equiv 0 \pmod 4$ and $x^t G x \le 2B$ with $G$ the Gram
matrix of twice the norm, which is $M$ odd and $M \le B$.  The box
formula agrees with line~6 (\code{quat\_\allowbreak{}lattice\_\allowbreak{}bound\_\allowbreak{}parallelogram},
\code{quaternion/\allowbreak{}ref/\allowbreak{}lvlx/\allowbreak{}lll\_\allowbreak{}applications.c}).  Line~5 reuses the
loop index $i$ of line~4.

\paragraph{Recommendation.}
State the skip on $\beta_i=0$ and the interval passed to the sampler,
using the corrected distribution in \S\ref{sec:sampler}. Reconcile the
loop policy with the failure analysis in \S9.5: specify the cap,
restart or error behavior, and its effect on the accepted distribution
and timing. Removing the cap or introducing retries requires checking
the assumptions of that analysis.

\section{Checked curve decoding and verifier failures}
\label{sec:verify}
\interopfinding{\algo{3.24}; \S3.1.2, \S3.4.1, \S3.4.3}{\medium}

The reference rejects $A_{\mathrm{pk}}=\pm2$ and
$A_{\mathrm{aux}}=\pm2$ before reconstructing bases~\cite{cref-verify}.
The specification already requires $A^2\ne4$ for \code{ECCurve} in
\S3.1.2. The useful clarification is where decoding untrusted bytes
establishes this invariant before routines that assume a valid curve
are called; the missing line in \textsc{Verify} is not by itself
evidence that the specification permits singular inputs.

Supporting calls on lines~5--11 may fail on malformed inputs.
Algorithm~3.24 should explicitly route these failures through the
exception convention of \S3.4.1 to \code{false}. A caller should not
need to infer whether a failed basis reconstruction, invalid kernel,
or unsuccessful split is an API error or an invalid signature.

The supersingularity claim in \S2.3.3 would also benefit from a short
argument linking the order and kernel checks to the properties needed
of the public and auxiliary curves. Section~3.4.3 expressly permits
deferred validation; \S\ref{sec:chainchecks} asks for that dependency
to be made explicit.

\paragraph{Suggested change.}
Show checked decoding of both coefficients before basis reconstruction,
state that supporting-routine failures reject the signature, and
identify the checks that justify each required curve and torsion
property. Include negative vectors for singular coefficients and
failed subroutines. No singular-curve acceptance or odd-response
verification bypass has been established here.

\section{Make deferred kernel checks traceable}
\label{sec:chainchecks}
\interopfinding{\algo{4.116}, \algo{4.118}, \algo{4.122},
  \algo{4.124}; \S3.4.3}{\medium}

The reference gates some kernel validation on its \code{verify} flag.
In particular, the gluing code explicitly verifies the projected
two-torsion points and their independence through
\code{verify\_two\_torsion}~\cite{cref-gluing}. The printed
\textsc{GluingCodomain} does not expose the same standalone check.
Related constraints also occur in supporting routines such as
\textsc{CoSymElem}, so the absence of that line alone does not establish
a difference in the accepted signatures.

Section~3.4.3 allows checks to be deferred to operations that enforce
the required property. What is missing for an implementor is a direct
mapping from each required order, independence, and splitting condition
to the operation that checks it, together with the treatment of failure.
Similarly, \algo{4.116} states its zero-coordinate error without the
reference's verification-mode qualification.

\paragraph{Suggested change.}
List the verifier's required predicates and cite the routine enforcing
each one. Distinguish mandatory validation of attacker-provided data
from signing-time consistency checks whose preconditions are already
established. Add negative vectors that reach each rejection point.
No signature accepted by the complete specified verifier but rejected
by the reference has been demonstrated for this concern.

\section{Ideal-to-isogeny kernel layout}
\label{sec:id2iso}

\interopfinding{Affects \algo{4.150} (\textsc{IdealToIsogeny}) lines~5 to~7
  and~13, and \algo{4.149} line~19.}{\medium}

The pseudocode places $(P_0, Q_0)$ on the first factor of
$E_0 \times E_0$, the images $\theta(P_0), c_2\theta(Q_0)$ scaled by
$\nu_1$ and $c_2 = \nu_1(1 - \nu_2 2^{f-2})$ on the second factor,
with $\theta = \beta_2 \overline{\beta_1}/N$, and reads $E_I$ from the
first codomain factor.  This matches Eq.~(28) in Appendix~A.

The reference does something else.  \code{quat\_\allowbreak{}qlapoty} returns
$\overline{\theta} = \beta_1 \overline{\beta_2}/N$ (it conjugates after
the norm equation).  \code{dim2id2iso\_\allowbreak{}ideal\_\allowbreak{}to\_\allowbreak{}isogeny\_\allowbreak{}qlapoty}
then forms
$T_1 = ([\nu_1]\overline{\theta}(P_0), P_0)$ and
$T_2 = ([\nu_1]\overline{\theta}(Q_0), Q_0)$, with no $c_2$ factor, runs
a dedicated chain variant (\code{theta\_\allowbreak{}chain\_\allowbreak{}compute\_\allowbreak{}and\_\allowbreak{}eval\_\allowbreak{}E2}),
and reads $E_I$ and the pushed basis from the \emph{second} codomain
factor.

Both recipes define the same $(2^{f-2}, 2^{f-2})$-isogeny kernel
$\langle [4]T_1, [4]T_2 \rangle$, and the outputs are $x$-only, so they
may well produce identical bytes.  We have not verified that a literal
\algo{4.150} does, and the pair $(T_1, T_2)$ the reference uses is not
isotropic at level $2^{f}$ in the sense required by \S4.5.4.1 (the
Weil pairing exponent is $\nu_1 2^{f-2}$), so it relies on chain
internals the spec does not describe.

\paragraph{Recommendation.}
State that \algo{4.150} was validated against the reference KATs, or
document the reference's kernel layout and the chain variant it needs.

\section{Parity loss in ideal intersection}
\label{sec:intersect}

\interopfinding{Affects \algo{4.77} (\textsc{Intersect}) via \algo{3.23}
  line~14 and \algo{3.21} line~4.}{\medium}

Line~2 reduces $X_2 + \mathbf{i}Y_2$ with \textsc{ModO$_0$}$(\cdot, N_2)$,
which maps each coordinate to $[0, N_2)$.  Lines~4 and~5 then divide
$X_2$ and $Y_2 - u^2 - v^2$ by two and assert the results are
integers.  Reduction modulo the odd $N_2$ does not preserve parity, so
the claimed integrality can fail after the reduction.
The reference (\code{quat\_\allowbreak{}ideal\_\allowbreak{}intersect\_\allowbreak{}O0}) reduces both
coordinates modulo $2N_2$, which keeps the parities of the unreduced
product.

\paragraph{Recommendation.}
Write line~2 with modulus $2N_2$, or reduce after the halving.

\section{Prime-norm sampling order and square-root choice}
\label{sec:primenorm}

\interopfinding{Affects \algo{4.71} (\textsc{PrimeNormIdealInertSampling})
  via \algo{4.81}, i.e.\ $I_{\mathrm{sk}}$ and
  $I_{\mathrm{com}}$.}{\medium}

Three details fix the bytes and are not in the text.
\begin{enumerate}[nosep]
  \item Draw order.  The reference
    (\code{quat\_\allowbreak{}random\_\allowbreak{}ideal\_\allowbreak{}O0\_\allowbreak{}given\_\allowbreak{}prime\_\allowbreak{}norm}) draws $x$,
    loops until the residuosity test passes, computes the square root,
    and only then draws the sign bit.  The pseudocode agrees, but
    nothing says the sign bit is a separate call to the interval
    sampler on $[0, 1]$, which costs nine bytes (\S\ref{sec:sample}).
  \item Root choice.  Line~5 calls \textsc{ModularSQRT}, whose output
    for $N \equiv 1 \pmod 4$ depends on the non-residue used by
    Tonelli--Shanks (\algo{4.20} does not fix it).  The sign bit of
    line~9 does not absorb this: the two roots differ in $Y$, hence in
    $y$, before and after the sign flip.
  \item The test in line~3 is implemented as
    $\left(\tfrac{-4x^2 - p \bmod N}{N}\right) = 1$ together with
    $(-4x^2 - p) \bmod 4 \in \{0, 1\}$.  This is equivalent, but the
    spec should say so since ``square mod $4N$'' has no algorithm
    attached.
\end{enumerate}

\paragraph{Recommendation.}
Pin the root returned by \textsc{ModularSQRT} (e.g.\ the smaller
representative, or the output of a stated Tonelli--Shanks
non-residue search), and state the byte cost of the sign draw.

\section{Primality-test witnesses and iteration counts}
\label{sec:isprime}
\interopfinding{\algo{4.23} and callers \algo{4.46}, \algo{4.79},
  \algo{4.148}}{\medium}

Line~18 samples a witness at random without naming a PRNG. The values of
\code{reps} (64, 96, 128) appear in the prose of \S4.1.3 but not in
Tables~8--10. At level~I the reference passes
\code{QUAT\_primality\_num\_iter} $=64$ and draws witnesses from the
default domain
(\path{precomp/ref/p324_3/include/quaternion_constants.h}).

Drawing witnesses from the \code{key}, \code{com}, or \code{res} stream
instead changes the later draws in that stream. Isolating the witnesses
in $\mathrm{prng}_{\mathrm{def}}$ avoids this source of interference.
The primality procedure is a probable-prime test: different witness
choices can affect acceptance of composites. It should not be described
as a Las Vegas computation whose result is independent of randomness.

\paragraph{Suggested change.}
Tabulate \code{reps}, name $\mathrm{prng}_{\mathrm{def}}$ in the algorithm,
and state the false-acceptance bound and assumptions used for this
probable-prime test. Keep its error accounting distinct from the
random-stream convention needed for conformance.

\section{Key generation and signing failure handling}
\label{sec:retry}

\interopfinding{Affects \algo{3.22} (\textsc{KeyGen}) and \algo{3.23}
  (\textsc{Sign}).}{\low}

The reference key generation (\code{protocols\_\allowbreak{}keygen},
\code{signature/\allowbreak{}ref/\allowbreak{}lvlx/\allowbreak{}keygen.c}) loops the whole of lines~3 and~4
until \textsc{IdealToIsogeny} succeeds, drawing fresh
\code{"key"}-domain bytes each time.  Signing
(\code{protocols\_\allowbreak{}sign}) returns failure on the two coprimality
exceptions of \algo{3.18} and \algo{3.21} and on
\textsc{RandomIdealGivenNorm} failure, and aborts on
\textsc{IdealToIsogenyNormEquation} failure.  The spec raises and
stops. Chapter~9 analyzes rare failures under its sampling assumptions,
but a small failure probability does not specify the API contract.
The distribution correction in \S\ref{sec:sampler} and the loop policy
in \S\ref{sec:rbe} should be reconciled with those assumptions.

\paragraph{Recommendation.}
State the intended behavior on each exception.

\section{Theta-to-Montgomery candidates and point maps}
\label{sec:theta2mont}

\interopfinding{Affects \algo{4.126} (\textsc{ThetaToMontgomery}) and
  \algo{4.96}.}{\low}

The third theta null point enumerated is $(a + \mathbf{i}b : a -
\mathbf{i}b)$ in the spec and $(\mathbf{i}a + b : a + \mathbf{i}b)$ in
the reference (\code{ec\_\allowbreak{}theta\_\allowbreak{}to\_\allowbreak{}montgomery},
\code{ec/\allowbreak{}ref/\allowbreak{}lvlx/\allowbreak{}normalize.c}).  The two give opposite Montgomery
coefficients, so after the $\pm A_i$ step of line~13 the same six
values are compared and the same $A_{\max}$ is chosen.  The conversion
matrices $\mathbf{N}$ of lines~20 to~31 are indexed by the position of
the maximum, and the positions of the third pair are swapped between
the two lists.  We did not validate the printed $\mathbf{N}$ table
against the reference's.

\paragraph{Recommendation.}
Either enumerate the same six null points as the reference or state
that any enumeration works and give the $\mathbf{N}$ matrices as
functions of the chosen null point rather than of an index.

\section{The v3 square-root convention}
\label{sec:sqrt}

\interopfinding{Affects \algo{4.35} (\textsc{Fp2SQRT}) and every consumer of
  \textsc{ProjectiveDifference}.}{\low}

\algo{4.35} matches the reference's \code{fp2\_\allowbreak{}sqrt}
(\code{gf/\allowbreak{}ref/\allowbreak{}lvlx/\allowbreak{}fp2.c}) line for line, including the $x_1 = 0$
select.  The v2 reference additionally negated the result to make the
real part even; that step is gone and \algo{4.35} is the whole
convention.  Nothing in \S1.4 says so.  An implementor porting a v2
field layer keeps the parity step, and every torsion basis differs.

\paragraph{Recommendation.}
Note in \S4.2.1.3 that the output of \algo{4.35} is not further
canonicalized, and list the change in \S1.4.

\part{Further specification clarifications}
\label{part:clarifications}

\section{Carry the difference point through basis actions}
\label{sec:actonbasis}

\interopfinding{Affects \algo{3.7} (\textsc{ActOnBasis}) and its callers
  \algo{4.150} lines~7 and~9, \algo{3.21} line~8, \algo{3.24}
  line~11.}{\medium}

\algo{3.7} returns the pair $([x_1]P + [x_3]Q,\, [x_2]P + [x_4]Q)$.
\algo{4.150} line~9 destructures a triple $(P_1, Q_1, R_1)$ from it
and pushes $R_1$ through the chain as the difference of the first two.
In $x$-only arithmetic the third point is a separate biscalar
multiplication by $(x_1 - x_2, x_3 - x_4)$ on the same triple, and it
cannot be recovered from the first two outputs (v2 review, item~17).
The reference's \code{id2iso\_\allowbreak{}matrix\_\allowbreak{}application\_\allowbreak{}even\_\allowbreak{}basis}
computes it when asked.

\paragraph{Recommendation.}
Define \textsc{ActOnBasis} on triples and state the third scalar pair.
Restate the rule the v2 review asked for, which v3 still leaves
implicit: once a basis $(P, Q, P - Q)$ exists, the third point is
carried through every doubling, matrix action, and isogeny evaluation
alongside $P$ and $Q$.  Recomputing it with
\textsc{ProjectiveDifference} picks a square-root branch independently
and can return $P+Q$ in place of $P-Q$, invalidating later
differential additions that rely on the original sign relation.

\section{Montgomery normalization and point transport}
\label{sec:normmont}

\interopfinding{Affects \algo{4.96} (\textsc{NormalizeMontgomery}) and
  \algo{4.99} (\textsc{TwoIsogenyChain}) line~26.}{\medium}

\begin{enumerate}[nosep]
  \item Line~1 feeds $(X_P + Z_P, X_P - Z_P)$ of a 4-torsion point to
    \textsc{ComputeMontCoeff}, which expects a one-dimensional theta
    null point.  The relation is Eq.~(10) of \S4.5.7.1 with
    $P = (a + b : a - b)$, thirty pages later, and is not cited.
  \item \textsc{ComputeMontCoeff} raises when $C = 0$.  \algo{4.96}
    does not say what happens then.  The reference
    (\code{ec\_\allowbreak{}theta\_\allowbreak{}to\_\allowbreak{}montgomery}) returns failure, which
    \algo{3.24} must treat as rejection.
  \item Changing the Montgomery model moves every point.  \algo{4.126}
    returns the $2 \times 2$ map $\mathbf{N}$; \algo{4.96} returns
    only the curve.  \algo{4.99} evaluates \code{pts} on line~24 and
    normalizes on line~26, so its output points are on the
    pre-normalization model.  No protocol call evaluates points
    through \textsc{TwoIsogenyChain}, so this is latent, but the
    algorithm as printed is inconsistent.  The reference's
    \code{ec\_\allowbreak{}normalize\_\allowbreak{}montgomery} applies the map to a basis when
    one is passed.
\end{enumerate}

\paragraph{Recommendation.}
Cite Eq.~(10) on line~1, add a failure clause, and return the point
map (or state that callers must not evaluate points).

\section{Work bounds for zero-hint basis reconstruction}
\label{sec:fromhint}
\interopfinding{\algo{4.95} lines~1--3 and the searches in \algo{4.94},
  reached from \algo{3.24}}{\low}

A zero hint makes \textsc{TorsionBasisFromHint} reconstruct the basis
by searching again. Algorithm~4.94 does not give a worst-case number
of candidates or a policy for search exhaustion. The reference also
uses search loops in \path{ec/ref/lvlx/basis.c}. This leaves the maximum
verification work and the treatment of an exhausted implementation
counter unclear.

This is a request for a work bound and a consistent failure policy.
We have not established an input that makes the v3 fallback search
nonterminating. Fast searches on sampled curves do not establish a
worst-case bound over the verifier's permitted input domain.

\paragraph{Suggested change.}
Give a justified search bound or a bounded reconstruction procedure,
and propagate its failure to signature rejection. Any cap that changes
the verifier's acceptance behavior needs an honest-input failure
argument and a common conformance rule. Specify counter exhaustion
explicitly, and test the last permitted candidate and failure path.
Do not infer a safe cap from an empirical maximum alone. The separate
$A=0$ basis convention is addressed in \S\ref{sec:hints}.

\section{Undefined batch normalization in change of basis}
\label{sec:normalize}

\interopfinding{Affects \algo{4.106} (\textsc{ChangeOfBasis}) line~2 and
  \algo{4.104} (\textsc{Tate}).}{\low}

Line~2 calls \textsc{Normalize} on ten points and the curve.  No
algorithm by that name exists.  \algo{4.104} requires its three points
as $(X : 1)$ and its curve as $(A_{24} : C_{24})$, and \algo{4.101}
divides by $C_{24}$ at every step.  The intent is a batch inversion
of the ten $Z$-coordinates (and presumably of $C_{24}$) with
\algo{4.36}.

\paragraph{Recommendation.}
Define \textsc{Normalize} or inline the batch inversion.

\section{Tie-breaking in Minkowski reduction}
\label{sec:minkowski}

\interopfinding{Affects \algo{4.65} (\textsc{MinkowskiReduce}) and the
  canonical-form claim of \algo{4.60}.}{\low}

Line~10 sorts the candidate list by norm.  Candidates of equal norm
have no stated order, and lines~13 to~25 take the first independent
ones.  \S9.1 argues that equal minima occur with probability far below
$2^{-\lambda}$, which covers the successive minima but not the
27-element candidate list, where several candidates share a norm
routinely (for instance $c$ and $-c$ are excluded by construction, but
$(1, 1, 0, 0)$ and $(1, -1, 0, 0)$ are not).  The output is claimed
canonical up to a global sign.

\paragraph{Recommendation.}
Specify a total order on candidates (norm, then the coefficient vector
lexicographically) or argue that the choice among equal-norm
candidates does not change $T$.

\section{The gluing ladder's output annotation}
\label{sec:gluingladder}

\interopfinding{Affects \algo{4.90} (\textsc{GluingLadder}).}{\low}

The output is annotated $[2^e - 1]P, [2^e]P$.  Starting from
$(R_0, R_1) = (Q, P)$ and applying \textsc{xDBLADD} $e$ times returns
$(R_1, R_0) = ([2^e]P,\; [2^e - 1]P + Q)$.  That is what \algo{4.131}
line~14 relies on, with $Q = 0$ on the first call and
$Q = [2^{m} - 1]T$ afterwards.

\paragraph{Recommendation.}
Annotate the output as $([2^e]P, [2^e - 1]P + Q)$.

\section{Incorrect fourth vector in the dual-lattice input}
\label{sec:dual4d}

\interopfinding{Affects \algo{4.60} (\textsc{DualLatReduction4DZ}).}{\low}

The input is written as
$2N\mathbb{Z} + 2N\mathbf{i}\mathbb{Z} + (A + B\mathbf{i} + \mathbf{j})\mathbb{Z} + (C + E\mathbf{i} + \mathbf{j})\mathbb{Z}$.
Doubling the basis of \algo{4.73} gives a fourth vector with
$\mathbf{k}$, not $\mathbf{j}$, and \algo{4.75} line~2 assigns
$C, E$ to the $\mathbf{k}$-column.  The Gram profile on line~2 of
\algo{4.60} is right for $\mathbf{k}$.

\section{Unit-normalization ties}
\label{sec:glzi}

\interopfinding{Affects \algo{4.59} (\textsc{LehmerGLZi}) lines~20 and~21,
  hence \algo{4.72} and \S\ref{sec:smallest}.}{\low}

Lines~20 and~21 pick $u \in \{\pm 1, \pm\mathbf{i}\}$ so that the
real and imaginary parts of $uA$ (respectively $uH$) are nonnegative.
When one part is zero two units qualify.  The choice fixes the first
reduced vector, which is $\alpha_1$ in \algo{4.80}, so it reaches the
key bytes.  Zero parts are rare but the rule should be total.

\section{A dual-isogeny mismatch in the Kani matrix}
\label{sec:kani}

\interopfinding{Affects \S4.6.1, Eq.~(15), and \algo{3.20}.}{\low}

The diagram above Eq.~(15) draws $\psi_2 : E' \to E_0$, but the
$(2,1)$ entry of $\Phi : E_0^2 \to E_I \times E'$ is written
$-\psi_2$, which would need domain $E_0$.  \S4.5.4 writes the same
entry as $-\psi_2$ with $\psi_2 : E_4 \to E_1$ and the same problem.
One of the arrow or the entry needs a hat.

\section{Secret-key decoding bounds}
\label{sec:skbounds}

\interopfinding{Affects Chapter~6 (secret key encoding).}{\low}

The inert ideal is encoded as three fields of
\code{FP\_\allowbreak{}ENCODED\_\allowbreak{}BYTES} bytes.  Nothing bounds
$N = \mathrm{IdealNorm}(I_{\mathrm{sk}})$ by $2^{328}$ at NIST-I; the
bound follows from \algo{4.79} and \S9.3 but is not stated at the
encoding.  No parse-time
checks are listed for a secret key ($N$ odd, $0 \le x, y < N$,
$M_{\mathrm{sk}}$ entries below $2^\lambda$ and normalized).  A
signer that accepts a tampered key runs \algo{3.23} on it.

\section{Editorial corrections}
\label{sec:editorial}

\interopfinding{Low-severity text defects.}{\low}

\begin{itemize}[nosep]
  \item Algorithms 4.42 to 4.48 are numbered out of order (4.43, 4.44,
    4.45, 4.42, 4.47, 4.46, 4.48).
  \item \algo{4.49} line~6 calls $\textsc{SumOfSquares}(r, m, \cdot)$;
    \algo{4.52} takes $(m, r, \beta_\lambda)$.
  \item \algo{4.79} line~9 calls $\textsc{InertGenToIdeal}(\beta, N')$;
    \algo{4.70} takes $(N, \gamma)$.
  \item \algo{4.80} line~1 omits the $\beta_\lambda$ argument of
    \algo{4.72}.
  \item \algo{4.83} line~5 reuses the loop index $i$ of line~4.
  \item \algo{3.21} line~12 writes $0_{E_{\mathrm{com}}}$ for a point
    of $E'_{\mathrm{aux}}$.
  \item \algo{3.13} is declared on \textsc{InertIdeal} with a footnote
    that it ``extends naturally''; \algo{4.83} line~3 applies it to a
    \textsc{MulIdeal}.  Say $N$ for all three types.
  \item Table~5 gives $e_{\mathrm{rsp}} = \lceil \log(p)/2 \rceil +
    \lambda/4 + 1$ and Table~8 gives
    $1 + \lceil \tfrac14(\lambda + 2\log_2 p) \rceil$.  They agree
    because $\lambda$ is a multiple of 8; one formula suffices.
  \item \algo{3.14} and \algo{3.15} write $\mathrm{cls}(()I)$.
  \item \S9.3.1: ``quadratif'', ``symetries du to'', ``we're gonna''.
  \item \algo{4.151} line~1 uses $M_{\mathbf{j}}$, which is not
    precomputed in \S4.6; give it as
    $2M_{(\mathbf{i}+\mathbf{j})/2} - M_{\mathbf{i}}$ or add it to
    \S4.7.
\end{itemize}

\part{Disposition of the v2.0.1 findings}
\label{sec:v2}

Each row cites the section of the v2.0.1 review~\cite{v2-review}.
\emph{Resolved} means v3 addresses the original issue, and \emph{moot}
means the relevant construction was replaced or removed.
\emph{Partial} marks an addressed issue with a related concern elsewhere
in v3. The zero-hint row is a request to clarify work bounds; it is not a
confirmed nontermination defect. These statuses concern the original
v2 finding, not every property of the associated v3 algorithm.

\begingroup
\small
\setlength{\tabcolsep}{4pt}
\begin{longtable}{@{}r >{\raggedright\arraybackslash}p{0.30\textwidth}
  >{\raggedright\arraybackslash}p{0.155\textwidth}
  >{\raggedright\arraybackslash}p{0.435\textwidth}@{}}
\toprule
\# & v2.0.1 finding & Status & v3.0 locus \\
\midrule
\endhead
1 & $\Fpii$ square root canonicalization & resolved &
  \algo{4.35} is a complete deterministic formula; the parity step is
  gone (\S\ref{sec:sqrt}). \\
2 & \textsc{ProductToTheta} coordinate convention & resolved &
  \algo{4.121} lines~4 to~8. \\
3 & ``Squared theta'' definition & resolved &
  \S4.5.2.1 separates \textsc{Squaring} and \textsc{Hadamard}; every
  algorithm spells both out. \\
4 & \textsc{GluingEval} cross-curve products & resolved &
  \algo{4.123} carries explicit subscripts (new formulas). \\
5 & Projective inverse of the dual null point & resolved &
  \algo{4.110} and \algo{4.122} lines~15 to~21 use products only. \\
6 & \code{hadamard\_\allowbreak{}bool} in the chain & resolved &
  \algo{4.132} passes $\mathrm{hadamard} = (m \ne 1)$ explicitly. \\
7 & Torsion basis slot convention & resolved &
  \S4.4.2 states the $(x_R, x_{RS}, x_S)$ permutation; \algo{4.94}
  line~31, \algo{4.95} line~14. \\
8 & Zero-hint fallback work bound & clarify &
  Bound and failure policy requested in \S\ref{sec:fromhint};
  no nonterminating v3 input established. \\
9 & Matrix storage convention & resolved &
  Eq.~(1), \algo{3.7}, Chapter~6 encoding order. \\
10 & $2^{f-e}$ scaling in \textsc{ChangeOfBasis} & resolved &
  \algo{3.8} defines $\mathcal{B}_2 = \mathcal{B}_1 \cdot 2^{f-e} M$
  with entries in $[0, 2^e)$; \algo{3.9} reduces first. \\
11 & Cubical Tate third-argument sign & resolved &
  Note after \algo{4.103}; \algo{4.106} inverts $\zeta_3, \zeta_4$
  explicitly. \\
12 & Matrix exponent omits \code{HD\_\allowbreak{}extra\_\allowbreak{}torsion} & resolved &
  \algo{3.23} line~18 uses $e_{\mathrm{rsp}} + 2$. \\
13 & Normalization of $(A_{24} : C_{24})$ & resolved &
  \algo{4.94}, \algo{4.95} require $(A : 1)$; \algo{4.99} ends in
  \textsc{NormalizeMontgomery}. \\
14 & Gluing requires Jacobian $y$ & moot &
  Superglue gluing is $x$-only (\algo{4.122} to \algo{4.124}). \\
15 & Jacobian doubling for the strategy & moot &
  \algo{4.114} doubles theta points. \\
16 & \textsc{DifferencePoint} normalization factor & resolved &
  \algo{4.92} line~4 states
  $\gamma = C(\overline{C}\,\overline{X_P Z_Q - Z_P X_Q})^2$. \\
17 & Never recompute $P - Q$ & partial &
  \algo{4.150} returns a triple; \algo{3.21} pushes only two points
  (\S\ref{sec:hdresponse}). \\
18 & Change-of-basis application convention & resolved &
  Eq.~(1), \algo{3.7}. \\
19 & \textsc{RandomIdealGivenNorm} $\beta$ range inclusive & resolved &
  \algo{4.82} with \algo{4.24} closed intervals; byte cost now differs
  (\S\ref{sec:sample}). \\
20 & \textsc{GeneralizedRepresent\allowbreak{}Integer} ranges & moot &
  \algo{4.46} samples $z, t \ge 1$ in $\Oo$ with no isogeny condition;
  new divergences in \S\ref{sec:sample}. \\
21 & Response lattice product vs.\ intersection & resolved &
  \algo{3.23} lines~14 and~15, \algo{3.18}. \\
22 & \textsc{LatticeSampling} underspecified & partial &
  Dual box specified in \algo{4.83}; distribution correction in
  \S\ref{sec:sampler} and draw/failure policy in \S\ref{sec:rbe}. \\
23 & $L^2$ deep-insertion order & moot &
  Floating-point LLL removed (\S4.3.3). \\
24 & \textsc{FixedDegreeIsogeny} $u^{-1}$ & moot &
  Removed; the scaling appears as $\nu_1, c_2$ in \algo{4.150}. \\
25 & Fixed-precision HNF & moot &
  Inert representation with bounds (Table~10). \\
26 & Endomorphism image $P - Q$ and Okeya--Sakurai & moot &
  No Jacobian lift; \algo{4.150} line~9 pushes $R_1$. \\
27 & \textsc{GetIndexSplitting} malformed input & resolved &
  \algo{4.128} raises; \algo{4.133} returns \code{false}. \\
28 & Signing key caches the generator & moot &
  The key stores $(N, x, y)$. \\
29 & \textsc{IdealToIsogeny} line~6 matrix formula & moot &
  Klopine removed; see \S\ref{sec:id2iso}. \\
30 & $\beta$ lives in the reduced ideal & resolved &
  \algo{4.149} line~20 maps $\beta_1$ back through $\gamma$. \\
31 & \textsc{SuitableIdeals} $v_2(\gcd)$ & moot & Removed. \\
32 & \algo{8.47} implementation variants & resolved &
  One \textsc{Isogeny22Chain}; kernel order $2^{e+2}$ stated in
  \S4.5.4.1 and \algo{3.20}. \\
33 & \textsc{IdealToIsogeny} output basis alignment & resolved &
  \algo{4.150} pushes $[\nu_1]\beta_1(\mathcal{B}_0)$ (\S4.6.1). \\
34 & \textsc{IdealToIsogeny} must propagate $P - Q$ & resolved &
  \algo{4.150} output is a triple. \\
35 & Splitting randomization & moot &
  Splitting is deterministic: \algo{4.126} keeps the largest of the six
  Montgomery coefficients. The v2 randomized representative-selection
  procedure is no longer used. \\
36 & Multi-order $\beta$ transport & moot & Removed. \\
37 & Connecting-ideal norm encoding & moot & Removed. \\
38 & Action-matrix order index & moot &
  Only $\Oo$ (\algo{4.134}). \\
39 & Sign error in \algo{4.9} line~11 & moot &
  \algo{3.24} line~9 decodes the aux basis at
  $e_{\mathrm{rsp}} + 2$ directly. \\
40 & $\mathrm{chl}$ top-bit canonicality & resolved &
  $\lambda = 128$ fills 16 bytes; \algo{3.4}. \\
41 & Lower bound on $e'_{\mathrm{rsp}}$ & moot &
  Chain length is fixed at $e_{\mathrm{rsp}}$. \\
42 & Extra-torsion-free chain lower bound & moot &
  One chain variant, fixed lengths. \\
\bottomrule
\end{longtable}
\endgroup

Totals: 22 resolved, 17 moot, 2 partial, and 1 clarification request.

% =====================================================================

\clearpage
\section*{Checks and source material}
\phantomsection
\label{sec:checks}
\addcontentsline{toc}{section}{Checks and source material}

The exact arithmetic checks accompanying the security findings run with
the Python standard library:
\begin{Verbatim}[fontsize=\small,frame=single,framesep=3mm]
python3 scripts/audit_sqisign_v3_design.py
\end{Verbatim}
They cover interval widths 1--512 at three draw widths, the three
\code{EIBox} calculations, the illustrative lattice rejection,
degree-9 kernel enumeration, degree-weight comparisons, and the
small-field matrix count. They establish the stated finite counts;
they do not simulate $\Hsim$ or run an end-to-end signer or verifier.

The conformance sections retain the earlier implementation comparisons
and their stated limits. This update corrects the affected sampler,
hint, and verifier recommendations; it does not rerun all reference
KATs or independently revalidate every inherited comparison.

The reviewed PDF is v3.0, dated 1 September 2026. Its SHA-256 digest is
\begin{quote}\small
\path{843912e5eb3dee20288af9fc9f01c0e1f1055834d3ce87c4687795c13adfd73c}.
\end{quote}
Reference paths in the findings are relative to \path{src/} at commit
\href{https://github.com/SQISign/the-sqisign/tree/6d017708db403bf83977fa70770fc4f7f9e9ff21}{\code{6d017708db40}}.
Ideal and quaternion norms are reduced norms; basis matrices act
column by column as in Eq.~(1) of the specification.

For tracking alongside the working notes, \emph{Biased interval sampling}
has the link target \hyperlink{SQI3-SPEC-01}{\code{SQI3-SPEC-01}},
and \emph{Hint degree weighting} has
\hyperlink{SQI3-PROOF-02}{\code{SQI3-PROOF-02}}.

\phantomsection
\addcontentsline{toc}{section}{References}
\begin{thebibliography}{9}
\small\raggedright

\bibitem{sqisign-spec}
SQIsign team.
\href{https://sqisign.org/spec/sqisign-20260901.pdf}{\emph{SQIsign:
Algorithm specifications and supporting documentation}, version 3.0},
1 September 2026. Sampler p.~28; bounded response sampling p.~51;
basis hints p.~58; Figure~4 p.~104; stronger security and timing
pp.~110--111.

\bibitem{sqisign-cref}
SQIsign team.
\href{https://github.com/SQISign/the-sqisign/tree/6d017708db403bf83977fa70770fc4f7f9e9ff21}{\emph{the-sqisign}, round-3 reference implementation},
tag \code{nist-v3}, commit \code{6d017708db40}.

\bibitem{cref-sampler}
SQIsign team.
\href{https://github.com/SQISign/the-sqisign/blob/6d017708db403bf83977fa70770fc4f7f9e9ff21/src/mp/ref/lvlx/mp.c\#L1711}{Reference integer sampler},
\path{src/mp/ref/lvlx/mp.c}, \code{ibz\_rand\_interval\_with\_domain}.

\bibitem{proof}
Aardal, Basso, De Feo, Patranabis, and Wesolowski.
\href{https://eprint.iacr.org/2025/379}{\emph{A Complete Security Proof of SQIsign}},
IACR ePrint 2025/379, consulted version dated 1 August 2025, Section~5.6.

\bibitem{cref-verify}
SQIsign team.
\href{https://github.com/SQISign/the-sqisign/blob/6d017708db403bf83977fa70770fc4f7f9e9ff21/src/verification/ref/lvlx/verify.c}{Reference verifier},
\path{src/verification/ref/lvlx/verify.c}, \code{protocols\_verify}.

\bibitem{cref-gluing}
SQIsign team.
\href{https://github.com/SQISign/the-sqisign/blob/6d017708db403bf83977fa70770fc4f7f9e9ff21/src/hd/ref/lvlx/gluing.c}{Reference gluing validation},
\path{src/hd/ref/lvlx/gluing.c}, \code{verify\_two\_torsion} and
\code{gluing\_compute}.

\bibitem{v2-review}
D.~Connolly.
\href{https://github.com/selkie-cryptography/sqisign-selkie/blob/16d7299bd096771ae5661b9fc21c5c98bbda39fd/latex/sqisign-v2-spec-review.tex}{\emph{SQIsign v2.0.1 Specification Review}},
Selkie Cryptography, 2026; archived source.

\end{thebibliography}
\end{document}
